\documentclass[11pt]{article}
\usepackage{preprint}
\usepackage{xeCJK}
\setCJKmainfont{FandolSong-Regular.otf}
\xeCJKDeclareCharClass{Default}{"2018,"2019,"201C,"201D}
\xeCJKsetup{CJKecglue={}}

\title{Rethinking the Gaokao: Charting a New Course for Education in China}
\author{Jiadi He\email{jh2288@georgetown.edu}}
\shorttitle{Rethinking the Gaokao}
\keywords{assessment; education policy; opportunity}

\begin{document}
\maketitle
\keywordsline

Within the educational sphere of China, the Gaokao commands significant attention. This crucial college entrance examination is deeply embedded in the fabric of Chinese society, serving as the primary barometer for academic success and future prospects. However, as China rapidly modernizes, the Gaokao's traditional focus on memorization increasingly misaligns with the broader range of skills necessary for today's dynamic global landscape.

The Gaokao's current model, which heavily leans on rote memorization, provides uniform scholastic benchmarks but does not adequately embrace the wide array of abilities needed in the modern workforce. Haiyan Zhang and Yubing Song's scholarly work brings to light the urgent necessity for an educational system that assesses a broader spectrum of student abilities beyond what a single examination can gauge, indicating a shift towards more versatile and comprehensive forms of assessment.

Reforms in China's education must thus tackle two pivotal concerns: expanding the scope of skills assessed and enhancing student well-being. The Gaokao's traditional structure, with its limited academic focus, falls short of capturing a student's full range of talents and their applicability to the professional world. Consequently, there is an increasing advocacy for incorporating diverse methods of assessment that recognize and value problem-solving, innovative thinking, and creativity.

In response, Chinese educational authorities are called upon to contemplate the adoption of a multi-faceted assessment approach. This could involve complementing the Gaokao with additional evaluations that better reflect students' abilities to innovate and apply their learning in real-world scenarios. Such an approach would not only uphold the stringent standards of the current system but also encourage an educational climate that fosters essential skills for the 21st-century success.

Additionally, the intense pressure associated with the Gaokao has far-reaching implications for students' mental health, with high levels of stress and anxiety being commonplace. The quest for academic distinction must balance with the imperative of student welfare. A shift towards a more rounded evaluation system, potentially incorporating ongoing assessments and project-based learning, is crucial in supporting a more nurturing educational experience.

The potential advantages of revising the Gaokao are extensive. For students, a more equitable assessment approach could mitigate the intense pressures they currently endure, promoting a healthier educational atmosphere. For educators, it would enable the exploration of new teaching methodologies aligned with pedagogical best practices, leading to a more engaging and diverse classroom experience. For universities, it could result in a richer pool of candidates assessed across a wider range of achievements and competencies.

Achieving these goals demands a collective and determined effort from all parties involved in China's education sector. It requires redefining educational success, moving from an exclusive emphasis on exam scores to a broader recognition of students' abilities and potential.

As China positions itself as a leader in global innovation, its education system stands at a transformative junction. The Gaokao, as the key to higher education, must adapt to mirror the values of a society that not only values academic rigor but also advocates for the comprehensive growth of its younger generation. By reshaping the Gaokao, China seizes the chance to establish a new paradigm for educational assessment, one that is as varied and complex as the students it aims to serve.

In sum, the reform of the Gaokao is more than an educational policy change; it signifies China's dedication to developing a generation capable of navigating the complexities of an ever-evolving global environment. The path ahead is substantial, yet the rewards for students, society, and China's global standing are potentially transformative. As China rethinks its educational aims, the Gaokao must also be transformed, not as an obstacle to be overcome but as a pathway to broader horizons and successes.

\begin{references}[Reference]
\item Zhang, Haiyan, and Yubing Song. On the Reform of Chinese High-entrance Exam from the Growing Tendency of the World’ s High-entrance Exam System, n.d.
\item “After 27 Tries, Chinese Millionaire Considers Giving up on College Entrance Exam.” NBCNews.com, June 28, 2023. \url{https://www.nbcnews.com/news/world/china-millionaire-gaokao-college-exam-27-tries-rcna91548?ex=digest}.
\item Yu, Yingshi. “谈中国高考制度的负面影响.” Radio Free Asia, October 11, 2020. \url{https://www.rfa.org/mandarin/pinglun/yuyingshi/yys-08022012121450.html}.
\item Gao, Yi. “广宇：中国高考制度的弊端 - 大纪元.” 大纪元 www.epochtimes.com, July 18, 2022. \url{https://www.epochtimes.com/gb/22/7/18/n13783352.htm}.
\item “专家：中英教育与高考制度利弊谈.” BBC News 中文. Accessed November 19, 2023. \url{https://www.bbc.com/zhongwen/simp/china/2013/06/130607_china_uk_edu}.
\end{references}

\end{document}
